%! AP Statistics — Unit 4, Lesson 4.2 — Warm-Up
\documentclass[10pt]{article}
\usepackage{apstats-article}
\usepackage{apstats-boxes}

\begin{document}

\pageheader{Unit 4, Lesson 4.2}{Warm-Up}
\namedateperiod

\vspace{0.3in}

% ── SECTION 1: Spiral Review ──────────────────────────────────────────────────
{\large\bfseries\color{navy} Part 1 \quad Spiral Review}

\vspace{0.12in}

\begin{enumerate}

    \item From Lesson 4.1 --- for each pattern, write the statistical question it suggests
    and state whether chance is a plausible explanation.

    \vspace{8pt}
    \renewcommand{\arraystretch}{1.8}
    \begin{tabularx}{\linewidth}{@{} X p{4cm} p{2.5cm} @{}}
    \textbf{Pattern} & \textbf{Statistical question} & \textbf{Chance plausible?} \\
    \hline
    A batter gets 6 hits in his last 7 at-bats (career average .280). &
    \blank{3.5cm} & \blank{2cm} \\
    A new drug trial shows 18 of 20 patients improved. &
    \blank{3.5cm} & \blank{2cm} \\
    \end{tabularx}

    \vspace{0.18in}

    \item A frequency table shows the results of 40 spins of a spinner with four equal
    sections (A, B, C, D):

    \vspace{6pt}
    \renewcommand{\arraystretch}{1.5}
    \begin{tabularx}{0.6\linewidth}{@{} l c c c c @{}}
    \textbf{Section} & A & B & C & D \\
    \hline
    \textbf{Count} & 8 & 12 & 11 & 9 \\
    \end{tabularx}

    \begin{enumerate}[label=(\alph*), itemsep=6pt]
        \item Compute the relative frequency for each section.\\[4pt]
        A: \blank{1.5cm} \quad B: \blank{1.5cm} \quad C: \blank{1.5cm} \quad D: \blank{1.5cm}
        \item If the spinner is fair, what relative frequency do you expect for each section?
              \blank{2cm}
        \item Does the data suggest the spinner is unfair? Explain.\\
        \writeline\writeline
    \end{enumerate}

\end{enumerate}

\vspace{0.35in}

% ── SECTION 2: Looking Ahead ──────────────────────────────────────────────────
{\large\bfseries\color{navy} Part 2 \quad Looking Ahead}

\vspace{0.12in}

\noindent Today we build a tool to estimate probabilities before we have formulas.

\vspace{0.14in}

\begin{tcolorbox}[colback=greenbg, colframe=greenacc, boxrule=0.8pt, arc=2mm,
    left=4mm, right=4mm, top=3mm, bottom=3mm]
    \small
    \begin{itemize}[leftmargin=1.3em, itemsep=8pt]
  \item A \textbf{simulation} uses a chance device to model a random process.
  \item We estimate probability by computing the \textbf{relative frequency} of the outcome over many trials.
  \item More trials $\to$ more accurate estimate (Law of Large Numbers).
    \end{itemize}
\end{tcolorbox}

\vspace{0.2in}

\noindent\textcolor{slate}{\small If you wanted to estimate the probability that a thumbtack
lands point-up when dropped, how would you go about it? Write a brief plan:}\\[8pt]
\writeline\\[2pt]\writeline

\end{document}
